% !TEX program = xelatex
\documentclass[12pt,a4paper]{article}
\usepackage{fontspec}
\usepackage[main=russian,english]{babel}
\usepackage[margin=25mm]{geometry}
\setmainfont{Times New Roman}
\newcommand{\ResultPending}[1]{\par\noindent\fbox{\parbox{0.94\linewidth}{Результат ожидается: #1}}\par}
\title{Шаблон новой публикации по диссертации}
\author{Авторские сведения заполняются после выбора площадки}
\date{Черновик}
\begin{document}
\maketitle
\begin{abstract}
Каркас предназначен для новой публикации. Он не изменяет поданную статью МИН-2026.
\end{abstract}
\section{Постановка задачи}
Симуляция, перенос на человеческие данные и клиническая SCS-проверка рассматриваются раздельно.
\section{Методы}
\ResultPending{зафиксировать эксперимент, baseline и независимое разбиение}
\section{Результаты}
\ResultPending{внести только воспроизводимые результаты после прохождения шлюза}
\section{Ограничения}
ECAP не является прямой мерой боли; симуляционная метка не является субъективной болью человека.
\end{document}
